% ====================================================================
%  v3_introduction.tex — Section 1: Introduction
%  Soccer-first framing per PAPER_PLAN_V2.md
%  Draft: v3 (2026-03-31)
% ====================================================================

\section{Introduction}\label{sec:intro}

Professional football operates under a soft-close rule.  A match has a
nominal duration of ninety minutes, but the referee adds discretionary
``stoppage time'' to compensate for interruptions---goal celebrations,
substitutions, injuries, VAR reviews.  Before 2022, second-half
stoppages averaged roughly four and a half minutes, despite ball-in-play
analyses suggesting that thirty to forty minutes per match were consumed
by non-playing events---a gap that referees systematically failed to
compensate, and in ways shown to reflect social pressure rather than
neutral judgment \citep{garicano2005favoritism, dohmen2008professionals}.
In November 2022, FIFA signaled to officials worldwide---first through
a directive to World Cup referees, then through league-level
adoption---that they should account more accurately for time lost.  The directive was not a new rule; Law~7 of the
Laws of the Game had always required full compensation.  It was an
\emph{enforcement shift}---a credible signal, demonstrated live during
the 2022 World Cup, that the gap between the law on the books and the
law in practice would close.

What happened to the game?  Second-half stoppage time in the five major
European leagues rose by 0.73 minutes relative to never-treated leagues
($p < 0.001$; wild cluster bootstrap $p = 0.036$), a 15\% increase over
the pre-policy mean.  The Premier League added the most (+1.31~min), the
Bundesliga and Ligue~1 roughly a minute each, La~Liga about two-thirds
of a minute; Serie~A, despite nominally adopting the directive, showed no
detectable change ($-0.09$~min, $p = 0.24$).  Including both halves,
total stoppage rose by 1.20 minutes per match.  Within treated leagues,
95--100\% of referees individually increased their stoppage time---near-universal
compliance achieved not through inspections, fines, or audits, but through
a demonstration event broadcast to two billion viewers.

The extra time changed the game's dramatic arc but not its aggregate
output.  Total goals per match were statistically unchanged (DiD $\beta =
-0.033$, $p = 0.39$).  Goals were \emph{displaced}, not created: scoring
in the second-half stoppage period rose by 35\% in treated leagues, while
the 76--84 minute window saw a commensurate decline.  Per-minute scoring
rates held constant across periods.  The quantity of football
increased; its per-minute intensity did not.  This is an 87\% quantity
effect: more minutes of play at unchanged rates, not a change in how the
game is played within each minute.

The displacement reshaped what matters most to viewers.  The share of
matches in which the result changed during stoppage time rose from 7.3\%
to 10.3\%---an additional ${\sim}55$ decisive moments per Big~5 season.  Late
equalizers increased by 23\%, late winners by 17\%.  A
\citet{gelbach2016when} decomposition reveals that 96--103\% of the
stoppage-time increase reflects referee discretion rather than
changes in match events.  Referees did not add time because more
stoppable events occurred; they added time because they were told to.
The downstream reallocation of goals, drama, and risk followed
mechanically.

\medskip

This pattern---more time, unchanged per-unit rates, displaced
outcomes---connects to a question that has occupied economists since
at least \citet{pencavel2015effects}: what happens when you extend the
duration of a high-intensity productive interaction?  Pencavel's
framework predicts concavity: beyond a threshold, additional hours yield
diminishing returns.  Our setting tests this prediction at an extreme
margin---minutes, not hours, at near-maximal physical intensity---and
finds approximate linearity.  The natural interpretation is that the
concavity in the hours-productivity relationship reflects cumulative
fatigue over long shifts; at very short extensions, fatigue does not
accumulate fast enough to bend the curve.  We cannot rule out that the
dose is simply too small to detect a declining slope, and we say so
explicitly.

The setting also speaks to \citet{rothockenfels2002}, who showed that
soft-close auction rules---where the deadline extends if activity
continues---largely eliminate deadline-driven strategic behaviour (``sniping'').
Professional football's stoppage rule is structurally parallel: the
contest has a nominal end, but the official extends it at discretion.
FIFA's directive increased the expected extension.  The Roth-Ockenfels
prediction is that strategic urgency near the deadline should diminish
when agents expect more time.  We find suggestive evidence consistent
with this: the sharp drop in goals during the 76--84 minute window
post-policy, combined with the reallocation into stoppage, is consistent
with teams pacing their late-match efforts differently when the
soft close extends.  But the mapping is not exact.  In
Roth-Ockenfels, bidders know the extension rule and optimize accordingly;
here, the regime change was unanticipated, and agents who built careers
under the old equilibrium had to \emph{learn} the new one.  This is
closer to real-world regulatory changes than the clean mechanism-design
setting of the original paper.

\medskip

Three features of the policy environment enable credible identification.
First, the directive created a well-documented, sharp treatment onset:
five leagues (Premier League, La~Liga, Ligue~1, Bundesliga, Serie~A)
adopted the enforcement shift for the 2023--24 season, while nine
leagues in our sample did not.  All five adopted within an eight-day
window (August 11--19, 2023); we are transparent that this is a
simultaneous shock, not meaningfully staggered adoption.  The design is
treated versus never-treated, with the eight-day wobble available as a
minor robustness check.  Second, the rich panel structure---15~leagues across
up to 13~seasons, over 36,000 matches with second-half stoppage
data---allows league and season fixed effects and informative
falsification tests.
Third, the directive targeted second-half stoppage specifically; the
first-half response ($\beta = +0.58$~min) provides a differential-dose
comparison, and pre-policy parallel trends (after excluding COVID-disrupted
seasons) strengthen the identifying assumption.

We complement the core difference-in-differences with three additional
identification strategies.  A sharp regression discontinuity at the
November 2022 World Cup start date exploits the within-league break in
enforcement practice.  Referee-level dose variation---407 officials with
heterogeneous pre-policy enforcement baselines---provides continuous
within-league identification.  And Callaway-Sant'Anna estimation,
appropriate for the treated-versus-never-treated structure, confirms
the main result ($\hat{\text{ATT}} = +0.49$~min, $p = 0.062$; pre-trends
$p = 0.90$).

\medskip

The paper traces the causal chain from mandate to on-pitch outcomes.
We find near-universal compliance through norm-based channels
(\S\ref{sec:compliance}), a large and robust first stage
(\S\ref{sec:firststage}), and a pure quantity effect with no
intensity change (\S\ref{sec:onpitch})---goals displaced into
stoppage time, not created; per-minute scoring rates unchanged; and
approximately 55~additional result-changing moments per Big~Five
season.

The paper contributes to several literatures.  To sports economics, we
provide the first large-scale causal estimate of the FIFA directive's
effects, resolving what had been an active debate about whether extended
stoppage time would increase total goals (it does not) or change the
game's competitive balance (it does, at the margin).  To labour
economics, we offer a test of the hours-productivity relationship at an
extreme margin that complements \citet{pencavel2015effects}'s
munitions-worker setting.  To the regulatory compliance literature, we
document a case of near-universal enforcement achieved through
demonstration rather than monitoring---a transmission mechanism
largely absent from the compliance frameworks of
\citet{johnson2020compliance} and \citet{levine2012randomized}.

\citet{garicano2005favoritism} showed that referees manipulate stoppage
time under social pressure; \citet{dugganlevitt2002} documented
corruption in competitive sports; \citet{pricewolfers2010} identified
racial bias in officiating.  Our contribution is distinct: we study
what happens when a global regulator \emph{succeeds} in changing
enforcer behaviour, and we trace the downstream consequences through
the game itself.

The remainder of the paper is organised as follows.
Section~\ref{sec:background} describes the institutional context and
identification strategy.  Section~\ref{sec:data} presents the data.
Section~\ref{sec:firststage} establishes the first stage.
Section~\ref{sec:compliance} examines referee compliance.
Section~\ref{sec:onpitch} decomposes the on-pitch effects.
Section~\ref{sec:conclusion} concludes.
